% Generated draft: requires mathematical and source review.
\documentclass[11pt]{article}
\usepackage[margin=25mm]{geometry}
\usepackage{fontspec}
\setmainfont{Latin Modern Roman}
\usepackage{amsmath,amssymb,mathtools,unicode-math}
\setmathfont{Latin Modern Math}
\usepackage{enumitem,needspace,xurl,hyperref}
\setlength{\parindent}{0pt}
\setlength{\parskip}{5pt}
\setlength{\emergencystretch}{3em}
\newcommand{\N}{\mathbb{N}}
\newcommand{\Z}{\mathbb{Z}}
\newcommand{\Q}{\mathbb{Q}}
\newcommand{\R}{\mathbb{R}}
\newcommand{\C}{\mathbb{C}}
\newcommand{\F}{\mathbb{F}}
\newcommand{\A}{\mathbb{A}}
\newcommand{\PP}{\mathbb{P}}
\newcommand{\E}{\mathbb{E}}
\newcommand{\eps}{\varepsilon}
\newcommand{\dd}{\,\mathrm{d}}
\begin{document}
\section*{UnsolvedMath ID: 700507}
\textbf{Source problem code:} catalogue entry
\subsection{Definitions and mathematical statement}
Let \(G\) be a topological group with identity \(1\). Say that \(G\) is locally nilpotent if every subgroup generated by finitely many elements is nilpotent, where a group \(H\) is nilpotent if its lower central series \(\gamma_1(H)=H\), \(\gamma_{k+1}(H)=[\gamma_k(H),H]\) reaches \(\{1\}\) after finitely many steps. Say that \(G\) is torsion-free if, for every \(g\in G\) and every integer \(n\geq 1\), \(g^n=1\) implies \(g=1\). Here the source's phrase ``without elements of finite order'' is understood to exclude nonidentity elements of finite order.

A subgroup \(N\leq G\) is isolated in \(G\) if
\[
\forall g\in G\;\forall n\in\Z_{\geq 1},\qquad g^n\in N\Longrightarrow g\in N.
\]
Closedness refers to the topology of \(G\), and normality means \(gNg^{-1}=N\) for every \(g\in G\).

The literal source question asks whether every locally compact, locally nilpotent, torsion-free, nonabelian topological group \(G\) has a subgroup \(N\) satisfying
\[
N\trianglelefteq G,\qquad N\text{ is closed},\qquad N\neq G,\qquad
\forall g\in G\;\forall n\in\Z_{\geq 1},\quad g^n\in N\Longrightarrow g\in N.
\]
The source does not require \(N\neq\{1\}\). Whether a nontrivial subgroup was intended is unresolved; that additional requirement has not been inserted. The source also does not explicitly specify a Hausdorff convention for locally compact groups.
\subsection{Short English statement}
Must every locally compact, locally nilpotent, torsion-free, nonabelian group contain a proper closed normal subgroup that contains every element having a positive power in that subgroup? The source leaves unclear whether the subgroup must also be nontrivial and whether Hausdorffness is assumed.
\subsection{Sources}
\begin{itemize}
\item Supplied local UnsolvedMath registry record: permanent ID \(700507\), exact source problem code \texttt{catalogue entry}, title ``Kourovka Notebook Problem 8.58''; upstream revision \texttt{d371ee917ac9647fbb06cd6fc529deec6f3dda82}.
\item Attributed to V. M. Poletskikh, proposed in 1982. The supplied background identifies \emph{The Kourovka Notebook (Unsolved Problems in Group Theory)}, 8th issue (1982), Problem 8.58, and the 21st issue edition (2026), \texttt{arXiv:1401.0300v48}. The collection description names E. I. Khukhro and V. D. Mazurov as editors.
\item The registry labels this problem ``open''. This is supplied source metadata, not an independently verified literature assessment. The cited Notebook text was not independently consulted for this draft.
\end{itemize}
\end{document}
